\begin{cabstract}
认知过程可以被视作具身化主体与环境持续交互中涌现的动态生成现象。脑电图（Electroencephalography, EEG）作为记录认知活动神经电信号的重要手段，在时空维度上直接反映了认知过程的动态生成特性。EEG通常来源于不同设备、任务范式与被试群体，形成了具有异质性的多源数据格局。
依赖特征工程与判别式的分析方法假设信号与神经模式之间存在相对稳定的映射关系，难以应对认知表征的可塑性与个体差异。
而生成式建模通过学习多源数据的潜在生成机制，能够从不完整观测中重构关键信息，从大规模无标注数据中学习可迁移的通用表征。
本文以多源EEG信号为研究对象，构建了从信号级到表征级的双层级生成建模体系，旨在实现多源EEG信号的深层分析。论文的主要工作及创新点包括以下方面：

（1）\textbf{多源EEG双层级生成式建模理论框架。}针对EEG信号的复杂时空结构和多尺度动态特性，本文构建了信号级和表征级的分层生成式建模架构，建立了从底层信号重构到高层语义挖掘的建模理论。信号级生成式建模方法学习EEG的内在生成机制，能够从低质量或不完整的观测中重构出高保真的神经活动模式，为后续分析提供可靠的数据。表征级生成式建模方法专注在无监督或弱监督条件下挖掘具有强泛化能力的潜在语义表征，从海量无标注数据中提取跨个体、跨任务的通用特征，有效缓解个体差异与域偏移带来的泛化困难。该框架通过形式化定义各层级的生成目标函数、结构化约束条件和层间信息传递机制，为多源EEG信号的理解提供了统一的理论基础。

（2）\textbf{基于状态空间建模的EEG超分辨率方法。}针对EEG空间分辨率不足的约束问题，本文提出了基于状态空间建模的脑电超分辨率方法（Mamba-based Super-Resolution for EEG, MASER）。作为信号级生成式建模中的数据增强策略，该方法实现了从稀疏电极配置推断高密度分布。首先，设计了基于状态空间模型的特征提取模块，有效捕获EEG信号的时序动态。其次，构建了低分辨率特征提取器和高分辨率信号预测器的级联生成架构，通过多尺度特征融合和渐进式上采样，实现了从粗粒度到细粒度的信号重建。第三，针对EEG信号的生理特性，设计了融合重建保真度损失、时域平滑性约束和连续性正则化的复合目标函数，确保生成信号在保持高保真度的同时具备良好的生理合理性。在SEED情感识别数据集和PhysioNet运动想象数据集上的实验验证表明，MASER在信号重建质量上显著优于现有的插值和深度学习基线方法，归一化均方误差降低16.25\%，皮尔逊相关系数提升1.13\%；在运动想象BCI任务中，通过4倍空间分辨率增强使分类识别准确率提高5.74\%，展现了超分辨率增强对下游应用任务的贡献。

（3）\textbf{协同引导学习驱动的EEG通道选择优化方法。}针对EEG设备在尺寸、重量、功耗和成本等方面的工程约束，本文提出了提出了协同引导学习驱动的认知状态空间模型（Cognitive State Space Model, CogSSM）。作为信号级生成式建模中的数据约简策略，该方法实现了在显著减少电极数量的前提下最大程度保持信号判别性能的智能优化。首先，受认知神经科学中全局工作空间理论和模块化专门处理机制的启发，设计了融合全局信息整合和任务特化处理的双路径协同引导学习架构。其次，开发了结合卷积和状态空间建模的核心特征提取模块，联合卷积操作和状态空间动力学建模局部时空模式与长程依赖关系。第三，构建了端到端的可微分通道重要性评估和选择机制，自动学习和发现对特定认知任务最具判别价值的电极位置组合。在OpenBMI和PhysioNet MI数据集上的验证显示，CogSSM仅使用20个通道即可达到76.95\%和88.55\%的准确率，相比基线方法提升了2.14\%-15.01\%，为便携式BCI系统的硬件优化提供了有效的解决方案。

（4）\textbf{基于去噪掩码自编码的EEG表征学习方法。}针对EEG数据标注成本高昂和跨域泛化困难的问题，本文构建了构建了脑电去噪掩码自编码器（Denoising Masked AutoEncoder for EEG, DMAE-EEG）预训练框架。作为表征级生成式建模的核心实现，该方法实现了从海量无标注数据中学习通用表征。首先，提出了基于脑区拓扑异质的数据划分方法，将均匀电极网格投影为反映大脑功能分区的层次化表征，有效解决了不同设备间电极配置差异问题。其次，设计了去噪伪标签生成器，通过重建型去噪任务为自监督学习提供稳健训练信号。第三，采用非对称掩码自编码器结构，通过轻量化解码器迫使编码器学习紧凑的语义表征。该框架在14个涵盖不同任务类型、设备配置和被试群体的公开数据集上进行预训练，通过零样本迁移和少样本微调验证了跨域迁移的有效性，在信号质量增强和运动意图识别任务的测试中均表现出显著的性能提升，体现了生成式自监督预训练在EEG信号分析中的巨大潜力。

基于上述双层级生成建模技术栈，本文设计并实现了认知状态监测与调节闭环系统原型。该原型构建了专注度--认知负荷二维评价体系，集成大语言模型驱动的建议生成、基于自适应音频的神经调制等功能模块，形成了感知、分析和干预的闭环管理，为生成式人工智能技术在脑机接口等领域的应用提供了参考。
\end{cabstract}

\ckeywords{脑电图；生成式建模；预训练；自监督学习；跨域迁移；掩码自编码器；脑机接口}

\begin{eabstract}
Cognitive processes can be viewed as dynamically emergent phenomena arising from the continuous interaction between embodied agents and their environments. Electroencephalography (EEG), as an important means of recording neural electrical signals of cognitive activity, directly reflects the dynamic generative characteristics of cognitive processes across spatiotemporal dimensions. EEG typically originates from different devices, task paradigms, and subject populations, forming a heterogeneous multi-source data landscape.
Analysis methods relying on feature engineering and discriminative approaches assume relatively stable mapping relationships between signals and neural patterns, making them inadequate for addressing the plasticity of cognitive representations and individual differences.
Generative modeling, by learning the latent generative mechanisms of multi-source data, can reconstruct key information from incomplete observations and learn transferable universal representations from large-scale unlabeled data.
We take multi-source EEG signals as the research object and construct a two-level generative modeling system spanning from signal level to representation level, aiming to achieve deep understanding and analysis of multi-source EEG signals. The main work and innovative points include the following aspects:

(1) \textbf{Multi-source EEG two-level generative modeling theoretical framework.} Targeting the complex spatiotemporal structure and multi-scale dynamic characteristics of EEG signals, we construct a hierarchical generative modeling architecture at signal and representation levels, establishing modeling theory from low-level signal reconstruction to high-level semantic mining. Signal-level generative modeling methods learn the intrinsic generative mechanisms of EEG and can reconstruct high-fidelity neural activity patterns from low-quality or incomplete observations, providing reliable data for subsequent analysis. Representation-level generative modeling methods focus on mining latent semantic representations with strong generalization capability under unsupervised or weakly supervised conditions, extracting universal features across individuals and tasks from massive unlabeled data, effectively alleviating generalization difficulties caused by individual differences and domain shift. The framework provides a unified theoretical foundation for understanding multi-source EEG signals through formal definition of generative objective functions at each level, structured constraint conditions, and inter-layer information transfer mechanisms.

(2) \textbf{EEG super-resolution method based on state space modeling.} Addressing the constraint problem of insufficient spatial resolution in EEG, we propose Mamba-based Super-Resolution for EEG (MASER), a method based on state space modeling for EEG super-resolution. As a data augmentation strategy in signal-level generative modeling, the method achieves inference of high-density distributions from sparse electrode configurations. First, a feature extraction module based on state space models is designed to effectively capture the temporal dynamics of EEG signals. Second, a cascaded generative architecture of low-resolution feature extractor and high-resolution signal predictor is constructed, achieving signal reconstruction from coarse-grained to fine-grained through multi-scale feature fusion and progressive upsampling. Third, targeting the physiological characteristics of EEG signals, a composite objective function integrating reconstruction fidelity loss, temporal smoothness constraints, and continuity regularization is designed to ensure that generated signals maintain high fidelity while possessing good physiological plausibility. Experimental validation on SEED emotion recognition dataset and PhysioNet motor imagery dataset shows that MASER significantly outperforms existing interpolation and deep learning baseline methods in signal reconstruction quality, reducing normalized mean square error by 16.25\% and improving Pearson correlation coefficient by 1.13\%; in motor imagery BCI tasks, 4× spatial resolution enhancement improves classification accuracy by 5.74\%, demonstrating the contribution of super-resolution enhancement to downstream application tasks.

(3) \textbf{EEG channel selection optimization method driven by collaborative guided learning.} Addressing engineering constraints of EEG devices in terms of size, weight, power consumption, and cost, we propose the Cognitive State Space Model (CogSSM) driven by collaborative guided learning. As a data reduction strategy in signal-level generative modeling, the method achieves intelligent optimization that maximally preserves signal discriminative performance while significantly reducing the number of electrodes. First, inspired by the global workspace theory and modular specialized processing mechanisms in cognitive neuroscience, a dual-path collaborative guided learning architecture integrating global information integration and task-specialized processing is designed. Second, a core feature extraction module combining convolution and state space modeling is developed, jointly modeling local spatiotemporal patterns and long-range dependencies through convolutional operations and state space dynamics. Third, an end-to-end differentiable channel importance assessment and selection mechanism is constructed to automatically learn and discover electrode position combinations most discriminative for specific cognitive tasks. Validation on OpenBMI and PhysioNet MI datasets shows that CogSSM achieves 76.95\% and 88.55\% accuracy using only 20 channels, improving by 2.14\%-15.01\% compared to baseline methods, providing an effective solution for hardware optimization of portable BCI systems.

(4) \textbf{EEG representation learning method based on denoising masked autoencoding.} Addressing the problems of high annotation cost and cross-domain generalization difficulties in EEG data, we construct the Denoising Masked AutoEncoder for EEG (DMAE-EEG) pretraining framework. As the core implementation of representation-level generative modeling, the method achieves learning of universal representations from massive unlabeled data. First, a data partitioning method based on brain region topological heterogeneity is proposed, projecting uniform electrode grids into hierarchical representations reflecting functional brain parcellation, effectively solving electrode configuration differences between different devices. Second, a denoising pseudo-label generator is designed to provide robust training signals for self-supervised learning through reconstruction-based denoising tasks. Third, an asymmetric masked autoencoder structure is adopted, forcing the encoder to learn compact semantic representations through a lightweight decoder. The framework is pretrained on 14 public datasets covering different task types, device configurations, and subject populations, and validates the effectiveness of cross-domain transfer through zero-shot transfer and few-shot fine-tuning, showing significant performance improvements in both signal quality enhancement and motor intention recognition tasks, demonstrating the great potential of generative self-supervised pretraining in EEG signal analysis.

Based on the above two-level generative modeling technology stack, we design and implement a prototype of a closed-loop system for cognitive state monitoring and regulation. The prototype constructs a two-dimensional evaluation system of attention and cognitive load, integrates functional modules such as large language model-driven recommendation generation and adaptive audio-based neuromodulation, forming a closed-loop management of perception, analysis, and intervention, providing reference for the application of generative artificial intelligence technology in fields such as brain-computer interfaces.
\end{eabstract}

\ekeywords{Electroencephalography, Generative modeling, Pretraining, Self-supervised learning, Cross-domain transfer, Masked autoencoder, Brain-computer interface}